\chapter{The maximal-slope identity}

\begin{lemma}[Slope bound from a metrized filtration]
\label{lem:filtered-slope-bound}\uses{lem:degree-calculus,lem:hn-filtration}\notready
If a saturated filtration of $A$ has semistable Hermitian quotients
of slopes at most $c$, then every nonzero submodule of $A$ has slope
at most $c$.
\end{lemma}
\begin{proof}
For a submodule $B$, intersect with every filtration step and take
successive quotients. Each quotient maps injectively and contractingly
into its actual image in the corresponding ambient grade: the ambient
quotient norm takes an infimum over a larger subspace. Semistability
bounds the degree of that image by $c$ times its rank. Contraction
bounds the quotient degree by the image degree. Add these inequalities;
quotient ranks sum to $\rk B$ and quotient degrees to $\degn B$.
Thus $\degn B\leq c\rk B$.
\end{proof}

\begin{theorem}[Maximal-slope identity]
\label{thm:bost-main}
\uses{lem:hn-filtration,lem:degree-calculus,thm:tensor-semistability,lem:filtered-slope-bound}
\source{bost-specht-manuscript:page-0001,bost-specht-manuscript:page-0013}
\notready
For every number field $K$ and every pair of positive-rank Hermitian
vector bundles $\overline E,\overline F$ over $\Ocal_K$,
\[
 \maxslope(\overline E\otimes\overline F)
 =\maxslope(\overline E)+\maxslope(\overline F).
\]
\end{theorem}
\begin{proof}
Let $0=E_0\subset\cdots\subset E_p=E$ and
$0=F_0\subset\cdots\subset F_q=F$ be their Harder--Narasimhan
filtrations. In lexicographic order use the steps
$E_{i-1}\otimes F+E_i\otimes F_j$ for $1\leq i\leq p$ and
$0\leq j\leq q$, identifying the repeated endpoint at each change of
$i$. Local splittings of the projective flag quotients show that this
is a saturated filtration whose grades are
$(E_i/E_{i-1})\otimes(F_j/F_{j-1})$. At each place take orthogonal
decompositions adapted to both flags. Their tensor decomposition is
orthogonal, so the grade quotient metrics are precisely the Hilbert
tensor metrics.

Each grade is semistable by tensor semistability, with slope the sum
of its factor slopes. This is at most
$c=\maxslope(E)+\maxslope(F)$. The filtration bound implies
$\maxslope(E\otimes F)\leq c$.
Conversely, choose saturated submodules $E'\subset E$, $F'\subset F$
attaining the two maximal slopes. Local splittings make
$E'\otimes F'\to E\otimes F$ injective, and the Hilbert metric is
the induced metric. Its average slope is $c$ by the tensor determinant
identity, so the maximal slope of the ambient tensor product is at
least $c$. The inequalities give equality.
\end{proof}

\begin{corollary}[Finite families]
\label{cor:bost-finite-family}\uses{thm:bost-main,lem:degree-calculus}
\source{bost-specht-manuscript:page-0001,bost-specht-manuscript:page-0013}
\notready
For every finite nonempty family of positive-rank Hermitian
$\Ocal_K$-bundles $(E_j)_{j=1}^N$,
\[
 \maxslope\left(\bigotimes_{j=1}^N E_j\right)
 =\sum_{j=1}^N\maxslope(E_j).
\]
If every factor is semistable, so is the tensor product.
\end{corollary}
\begin{proof}
Induct on $N$. The one-factor identity is tautological; the inductive
step applies the two-factor theorem to the tensor product of the
first $N-1$ factors and the last one. Hilbert tensor associativity is
an isometry, so parenthesization does not alter slopes. For semistable
factors the sum of their maximal slopes is the sum of their average
slopes, which is the average slope of the tensor product. Equality of
maximal and average slope is exactly semistability.
\end{proof}
